\scanpage{24}
% L1–2
\uline{Polyèdre combinatoire} associé à un polyèdre convexe :
l'ens des facettes \inserted{$\ne$ l'unique facette maximale}, \deleted{\unsure{autre que}} ordonné par $\overline F\subset\overline{F'}$.

% L3–8
[ La dim d'une facette se retrouve sur le
polyèdre combinatoire comme une dim.
combinatoire (d'un él d'un ens. ordonné) :
sup des longueurs des chaînes ascendantes
$F_0\lneq F_1\lneq\ldots\lneq F_n=F$ -- ou longueur commune
des chaînes maximales \ldots{} ]

% L9–13
\[
\unsure{P|\mathcal P|\Phi}_0\quad P_1\quad\text{--}\quad
P_{n-1}\deleted{\unsure{\{C\}|[C]}}\qquad (n=\dim C)
\]
\unsure{où|Si} $P_i=$ ens des facettes de dim. $i$
Entre $P_i$ et $P_{i+1}$ $(0\leq i\leq n-1)$ relation d'incidence \inserted{$R_i$}.
La relation d'ordre sur $P\simeq\coprod_{0\leq i\leq n}P_i$ se
réconstitue à partir des $R_i$.

\begin{center}-- -- -- -- --\end{center}

% L14–20. Left vertical rule.
\uline{Proposition} Les facettes \inserted{$\ne\varnothing$} du \deleted{\unsure{$(B,R)$}} \inserted{cube $C(B,R)$} sont en
corr. biunivoque avec l'ens des \deleted{\unsure{f}} Sections
partielles de $B$ sur $I=B/R$. Si
\diagram{\begin{tikzcd}[ampersand replacement=\&,column sep=small]
\sigma\colon J\arrow[rr,bend right=17] \&[-1.4em] \subset I \& B
\end{tikzcd}}
est une telle section, la
facette \uline{fermée} \inserted{(ouverte)} associée est l'ens des
$\displaystyle\sum_{b\in B}\lambda_b.b$ avec $\lambda_{-b}=-\lambda_b$, $\lambda_{\sigma(\unsure{j|y})}=+1$, et
$|\lambda_b|\leq 1$ (resp $|\lambda_b|<1$) pour $b\in I\unsure{-|\setminus}J$.
% archive sheet number: 23
